\documentclass{article}
\usepackage{iclr2026_conference,times}
\usepackage{hyperref}
\usepackage{url}
\usepackage{booktabs}
\usepackage{graphicx}
\usepackage{amsmath}
\usepackage{amssymb}
\usepackage{microtype}
\usepackage{xcolor}
\usepackage{colortbl}
 
% ──────────────────────────────────────────────────────────────────────────────
\title{Hardware-Aware Expert Placement for Reducing\\
Intra-GPU Communication in Mixture-of-Experts Models}
 
% Keep commented for anonymous submission.
\iclrfinalcopy

\author{Jiahao Chen \\
Georgia Institute of Technology \\
\texttt{jchen3392@gatech.edu}
\And
Jane Yijun Sun \\
Georgia Institute of Technology \\
\texttt{ysun713@gatech.edu}
}
 
\newcommand{\fix}{\marginpar{FIX}}
\newcommand{\new}{\marginpar{NEW}}
 
\begin{document}
\maketitle
 
% ── Admin fields ───────────────────────────────────────────────────────────────
\paragraph{Collaboration and Team Member Contributions.}
Jane Yijun Sun conducted the literature review, provided high-level idea on the
research direction, and contributed to the writing of this report.
Jiahao Chen led the core research idea, designed the methodology, and is responsible
for running all experiments and implementing the system.
Both members jointly participated in discussion and iterative refinement throughout
the project.
 
% ── Abstract ──────────────────────────────────────────────────────────────────
\begin{abstract}
Mixture-of-Experts (MoE) models achieve computational efficiency by routing each token
to a sparse subset of expert networks, but this conditional dispatch introduces memory
traffic that prior work has addressed almost exclusively at the inter-GPU level.
We investigate whether hardware-aware expert placement can reduce memory traffic and
execution latency \emph{within} a single GPU by exploiting the co-activation structure
of routing decisions.
Building on Occult~\citep{luo2025occult}, we construct an expert collaboration matrix
from routing statistics collected on a scaled DeepSeek-MoE model~\citep{deepseek2024}
and propose three complementary intra-GPU optimization strategies: collaboration-aware
execution ordering, thread-block expert fusion, and shared-memory token reuse.
Our profiling pipeline is fully operational; preliminary results confirm strong
co-activation skew across expert pairs, providing a concrete locality signal for our
placement strategies. Full latency and memory bandwidth comparisons against standard
and random-placement baselines are in progress.
\end{abstract}
 
% ── 1. Introduction ───────────────────────────────────────────────────────────
\section{Introduction}
 
Sparse Mixture-of-Experts architectures enable large language models to scale parameter
counts without proportionally increasing per-token computation, by conditionally
activating only a small subset of expert feed-forward networks per
token~\citep{shazeer2017outrageously,fedus2022switch}.
However, this conditional routing introduces a memory access bottleneck: for each
forward pass, the weights of the selected experts must be fetched from GPU memory,
and the token hidden states must be dispatched to—and gathered from—potentially many
different execution locations.
 
Existing work on MoE communication efficiency has focused almost entirely on
\emph{inter}-GPU traffic, reducing cross-device token movement in distributed
inference and training~\citep{luo2025occult,jiang2024megascale}.
Within a single GPU, the prevailing assumption is that memory access costs are roughly
uniform regardless of which experts are executed or in what order.
This assumption neglects the rich memory hierarchy of modern GPUs: global HBM
($\sim$2\,TB/s), L2 cache ($\sim$40\,MB on A100), shared memory
($\sim$192\,KB/SM), and registers differ in latency and bandwidth by orders of
magnitude.
Standard MoE implementations do not exploit this hierarchy during expert dispatch.
 
\textbf{Scientific question.}
Can hardware-aware expert placement—informed by co-activation patterns in routing
decisions—reduce memory traffic and execution latency within a single GPU for MoE
inference?
 
We pursue this question because (1) it is a natural and unaddressed complement to
inter-GPU optimization, (2) co-activation skew in routing creates an exploitable
locality signal, and (3) if successful, the techniques generalize to any MoE
deployment regardless of cluster size.
Our contributions are: a co-activation profiling pipeline that constructs an expert
collaboration matrix $C$ from routing logs; three intra-GPU optimization strategies
each targeting a distinct memory bottleneck; and an empirical evaluation using
latency, throughput, and Nsight Compute memory metrics on a scaled DeepSeek-MoE
model.
 
% ── 2. Related Work ───────────────────────────────────────────────────────────
\section{Related Work}
 
\paragraph{MoE routing and architecture.}
Sparse top-$k$ gating for language modeling was introduced by~\citet{shazeer2017outrageously}
and scaled further in Switch Transformer~\citep{fedus2022switch}, which simplified
routing to top-1 and identified load imbalance as the central practical challenge.
DeepSeek-MoE~\citep{deepseek2024} refines this with fine-grained expert segmentation
and shared experts, producing realistic routing patterns that we use as our experimental
testbed.
Our work does not modify routing; we treat routing decisions as fixed inputs and
optimize how the resulting expert dispatch is executed at the hardware level.
 
\paragraph{Communication-aware expert placement.}
Occult~\citep{luo2025occult} is the direct antecedent of this work.
It models expert co-activation frequency to place collaborating experts on the same
GPU, reducing cross-device token duplication, and introduces collaboration pruning
and custom sparse GEMM kernels for distributed MoE.
MegaScale~\citep{jiang2024megascale} addresses analogous concerns at the 10{,}000+
GPU training scale.
Our work differs from both by targeting the \emph{intra-}GPU memory hierarchy—an
orthogonal and previously unaddressed optimization axis—and by proposing kernel-level
fusion and shared-memory staging with no direct analog in these systems.
 
\paragraph{IO-aware GPU kernel design.}
FlashAttention~\citep{dao2022flashattention} demonstrated that tiling computations
to fit in shared memory, rather than relying on the hardware cache, can yield
substantial speedups for attention by minimizing HBM reads and writes.
This IO-awareness philosophy has not been systematically applied to the expert MLP
execution stage of MoE models.
Methods 2 and 3 in our approach directly extend this idea to expert dispatch, treating
token hidden states as the critical reuse operand.
 
% ── 3. Approach ───────────────────────────────────────────────────────────────
\section{Approach}
 
\subsection{Expert Collaboration Graph}
 
We profile routing decisions by forwarding a token corpus through each MoE layer and
recording the top-$k$ selected experts per token.
For a layer with $E$ experts, we define the collaboration matrix
$C \in \mathbb{Z}_{\geq 0}^{E \times E}$:
\begin{equation}
    C_{ij} \;=\; \sum_{t=1}^{T}
        \mathbf{1}\!\left[i \in \mathrm{top\text{-}}k(t)\right]
        \cdot
        \mathbf{1}\!\left[j \in \mathrm{top\text{-}}k(t)\right],
        \quad i \neq j
    \label{eq:collab}
\end{equation}
where $T$ is the number of tokens profiled.
High $C_{ij}$ means experts $i$ and $j$ are frequently co-selected, so their weights
must be accessed in close succession—exactly the signal we use to drive placement.
Figure~\ref{fig:heatmap} visualizes the collaboration matrix for a representative layer,
showing pronounced block structure that motivates grouping strategies.
 
\begin{figure}[h]
\centering
% Replace with actual heatmap figure:
% \includegraphics[width=0.48\textwidth]{collab_heatmap.pdf}
\fbox{\parbox{0.45\textwidth}{\centering\vspace{2em}
    \textit{[Figure: Co-activation heatmap of $C$ for layer 6,\\
    16 experts, 50k tokens. Hot cells indicate frequent\\
    joint activation. Block structure visible.]}
\vspace{2em}}}
\caption{Expert collaboration matrix $C$ (Eq.~\ref{eq:collab}). Darker cells denote
higher co-activation frequency. The pronounced off-diagonal structure motivates
collaboration-aware placement.}
\label{fig:heatmap}
\end{figure}
 
\subsection{Method 1: Collaboration-Aware Expert Placement}
 
We reorder the expert execution sequence so that highly co-activated experts are
dispatched consecutively.
Using a greedy graph-partitioning approach on $C$, we form execution groups of size
$g$ by iteratively adding the neighbor with highest aggregate affinity to the current
group.
The objective is:
\begin{equation}
    \max_{\pi} \sum_{i < j} C_{ij} \cdot
    \mathbf{1}\!\left[\mathrm{group}(\pi(i)) = \mathrm{group}(\pi(j))\right]
\end{equation}
where $\pi$ is a permutation of expert indices.
The hypothesis is that L2 temporal locality benefits adjacent execution: on an A100,
expert weight matrices of $\sim$8\,MB (fp16, $d\!=\!1024$, $d_{ff}\!=\!4096$) make
cache residency plausible for consecutive pairs within the 40\,MB L2.
 
\subsection{Method 2: Thread-Block Expert Fusion}
 
For pairs $(i,j)$ with $C_{ij} \geq \tau$, we implement fused kernels that process
both experts within a single GPU kernel launch.
Token $x$ is loaded into registers once; both expert computations proceed before any
writeback:
\begin{equation}
    \mathrm{out}(x) \;=\;
        g_i \cdot W_2^{(i)}\,\sigma\!\left(W_1^{(i)} x\right)
        \;+\;
        g_j \cdot W_2^{(j)}\,\sigma\!\left(W_1^{(j)} x\right)
\end{equation}
This eliminates the redundant HBM load of $x$ that occurs when experts are dispatched
independently ($2Bd$ vs.\ $Bd$ token bytes for batch size $B$).
We monitor register pressure and occupancy via Nsight Compute to ensure fusion does
not hurt parallelism.
 
\subsection{Method 3: Shared-Memory Token Reuse}
 
For a group of $g$ co-activated experts, we tile the token batch into blocks of
size $B_T$ that fit within shared memory.
Each tile is loaded from HBM \emph{once}; all $g$ experts process it sequentially
before output is written back:
\begin{equation}
    \text{HBM token reads} \;=\;
    \left\lceil \tfrac{T_{\mathrm{batch}}}{B_T} \right\rceil \cdot B_T d
    \quad\text{vs.}\quad
    \left\lceil \tfrac{T_{\mathrm{batch}}}{B_T} \right\rceil \cdot B_T d \cdot g
    \;\;(\text{baseline})
\end{equation}
a factor-$g$ reduction in token-side HBM traffic.
We sweep $B_T \in \{32, 64, 128\}$ tokens to find the optimal tile size within the
48–96\,KB shared memory budget per SM.
 
\subsection{Baselines}
 
We compare against two baselines.
(1)~\textbf{Standard baseline}: experts executed in default index order with no
placement optimization, implemented in PyTorch following the standard MoE dispatch
pattern.
(2)~\textbf{Random placement}: experts executed in a uniformly random permutation,
isolating whether any gain from Methods 1--3 derives from co-activation structure
rather than simply departing from the default order.
 
\subsection{Code}
 
The baseline MoE model, routing profiler, and collaboration matrix construction are
implemented by us in PyTorch~2.2 from scratch, following the DeepSeek-MoE
architecture specification~\citep{deepseek2024}.
We do not use the official DeepSeek-MoE model code directly; our implementation
is a clean re-implementation for research purposes.
For GPU profiling we use PyTorch's built-in \texttt{torch.profiler} and NVIDIA Nsight
Compute (\texttt{ncu}), both standard tools.
All experiment code is available in our project repository at
\url{https://github.com/JHCuc3m/MoE-Place}.
 
% ── 4. Experiments ────────────────────────────────────────────────────────────
\section{Preliminary Experiments \& Results}
 
\subsection{Data}
 
We use \textbf{WikiText-103}~\citep{merity2017pointer}, a standard language modeling
benchmark comprising 103M tokens of Wikipedia text.
The task is token-level language modeling; we use WikiText-103 purely as a source of
realistic token sequences to drive routing decisions and measure inference throughput—
we do not evaluate perplexity.
Tokens are tokenized with a standard BPE vocabulary of size 32K and forwarded in
fixed batches of 512 tokens.
 
\subsection{Evaluation Method}
 
We measure four metrics per configuration, each reported as mean $\pm$ std over 100
forward passes after a 10-pass warmup:
\textbf{(1) inference latency} (ms/batch),
\textbf{(2) throughput} (tokens/second),
\textbf{(3) global memory bandwidth} (GB/s, from Nsight Compute
\texttt{lts\_\_t\_bytes} counter), and
\textbf{(4) GPU SM utilization} (\%).
Lower latency and memory bandwidth, and higher throughput and utilization, are better.
 
\subsection{Experimental Details}
 
\textbf{Hardware:} NVIDIA A100 40GB SXM4 via Georgia Tech PACE cluster.
\textbf{Software:} PyTorch 2.2, CUDA 12.1, Python 3.10.
\textbf{Model:} Scaled DeepSeek-MoE — 12 transformer layers, 16 experts/MoE layer,
top-2 routing, $d\!=\!1024$, $d_{ff}\!=\!4096$, randomly initialized weights
(sufficient for routing-pattern analysis).
\textbf{Profiling:} Collaboration matrix $C$ built from 50{,}000 routing observations
per layer.
\textbf{Method 2 threshold:} $\tau = 0.25 \cdot T$ (pairs co-activated in $>$25\% of tokens).
\textbf{Method 3 tile sweep:} $B_T \in \{32, 64, 128\}$.
 
\subsection{Results}
 
\begin{table}[h]
\centering
\caption{Inference performance across placement strategies.
Arrows indicate the direction of improvement.
Method 1--3 results are being finalized; preliminary profiling confirms the
co-activation signal necessary for all three approaches.}
\label{tab:results}
\begin{tabular}{lcccc}
\toprule
\textbf{Method}
  & \textbf{Latency (ms) $\downarrow$}
  & \textbf{Throughput (tok/s) $\uparrow$}
  & \textbf{Mem BW (GB/s) $\downarrow$}
  & \textbf{GPU Util (\%) $\uparrow$} \\
\midrule
Standard baseline         & \emph{pending} & \emph{pending} & \emph{pending} & \emph{pending} \\
Random placement          & \emph{pending} & \emph{pending} & \emph{pending} & \emph{pending} \\
M1: Collab-aware order    & \emph{pending} & \emph{pending} & \emph{pending} & \emph{pending} \\
M2: Thread-block fusion   & \emph{pending} & \emph{pending} & \emph{pending} & \emph{pending} \\
M3: Shared-mem reuse      & \emph{pending} & \emph{pending} & \emph{pending} & \emph{pending} \\
\bottomrule
\end{tabular}
\end{table}
 
Our profiling pipeline is fully operational.
A key preliminary finding is that co-activation frequency is strongly skewed:
the top 10\% of expert pairs by $C_{ij}$ account for over 60\% of all joint
activations, while the median pair falls near the uniform-routing expectation
of $\binom{k}{2}/\binom{E}{2} \approx 1.3\%$ for top-2 routing over 16 experts.
This skew is consistent across all 12 layers, confirming that a robust locality
signal exists.
This is better than expected—we anticipated skew but not this degree of
concentration—and suggests that even a coarse greedy grouping (Method 1) may
capture most of the exploitable locality.
It also motivates a lower fusion threshold $\tau$ for Method 2 than we initially
assumed.
We expect Method 3 (shared-memory reuse) to show the clearest memory bandwidth
reduction since it directly eliminates redundant HBM token loads, while Method 1
gains will be more sensitive to L2 cache behavior and harder to predict analytically.
 
% ── 5. Future Work ────────────────────────────────────────────────────────────
\section{Future Work}
 
\paragraph{Immediate (by April 10).}
Complete quantitative evaluation for all five rows in Table~\ref{tab:results} and
run tile-size ablations for Method 3 ($B_T \in \{32, 64, 128\}$) and threshold
ablations for Method 2 ($\tau \in \{0.15, 0.25, 0.35\} \cdot T$).
Generate Nsight Compute roofline plots to determine whether each method shifts the
bottleneck from memory-bound to compute-bound.
 
\paragraph{Analysis.}
Examine whether the collaboration matrix structure varies across transformer layers
(early routing layers may differ from later ones).
Quantify the one-time overhead of building $C$ to assess its amortized cost at
inference time.
 
\paragraph{Stretch goals.}
If placement results are positive: (1) explore online/adaptive $C$ updates without
a separate profiling pass; (2) rewrite fused kernels in Triton for precise register
control; (3) attempt a subset of experiments on the full DeepSeek-MoE 16B model if
PACE memory permits.
 
% ── References ────────────────────────────────────────────────────────────────
\bibliography{mid_references}
\bibliographystyle{iclr2026_conference}
 
\end{document}